\documentclass[aspectratio=169,12pt]{beamer}
\usepackage[utf8]{inputenc}
\usepackage[T1]{fontenc}
\usepackage[spanish,mexico]{babel}
\usepackage{amsmath,amssymb,mathtools}
\usepackage{xcolor}
\usetheme{default}
\setbeamertemplate{navigation symbols}{}
\setbeamertemplate{footline}{%
  \leavevmode\hbox{%
  \begin{beamercolorbox}[wd=.82\paperwidth,ht=2.7ex,dp=1.2ex,leftskip=1em]{author in head/foot}%
    Cálculo Integral -- TecNM -- M. en C. Carlos E. Martínez R.
  \end{beamercolorbox}%
  \begin{beamercolorbox}[wd=.18\paperwidth,ht=2.7ex,dp=1.2ex,center]{date in head/foot}%
    \insertframenumber/\inserttotalframenumber
  \end{beamercolorbox}}%
}
\definecolor{azulcurso}{RGB}{25,67,120}
\setbeamercolor{frametitle}{fg=azulcurso}
\setbeamercolor{structure}{fg=azulcurso}
\setbeamerfont{frametitle}{series=\bfseries,size=\large}
\newcommand{\ejercicio}[1]{%
  \begin{center}
  \vfill
  {\fontsize{24}{30}\selectfont #1}
  \vfill
  \end{center}
}
\title{Integrales indefinidas}
\subtitle{Colección de ejercicios para trabajo en clase}
\author{M. en C. Carlos E. Martínez R.}\institute{Cálculo Integral \\ Tecnológico Nacional de México}
\date{}
\begin{document}
\begin{frame}[plain]
\titlepage
\end{frame}
\begin{frame}{Integración inmediata: regla de la potencia}
\small Ejercicio 1 \hfill Cálculo Integral -- TecNM
\ejercicio{$\displaystyle \int x^7\,dx$}
\end{frame}

\begin{frame}{Integración inmediata: regla de la potencia}
\small Ejercicio 2 \hfill Cálculo Integral -- TecNM
\ejercicio{$\displaystyle \int x^{12}\,dx$}
\end{frame}

\begin{frame}{Integración inmediata: regla de la potencia}
\small Ejercicio 3 \hfill Cálculo Integral -- TecNM
\ejercicio{$\displaystyle \int x^{-3}\,dx$}
\end{frame}

\begin{frame}{Integración inmediata: regla de la potencia}
\small Ejercicio 4 \hfill Cálculo Integral -- TecNM
\ejercicio{$\displaystyle \int x^{-5}\,dx$}
\end{frame}

\begin{frame}{Integración inmediata: regla de la potencia}
\small Ejercicio 5 \hfill Cálculo Integral -- TecNM
\ejercicio{$\displaystyle \int x^{1/2}\,dx$}
\end{frame}

\begin{frame}{Integración inmediata: regla de la potencia}
\small Ejercicio 6 \hfill Cálculo Integral -- TecNM
\ejercicio{$\displaystyle \int x^{3/2}\,dx$}
\end{frame}

\begin{frame}{Integración inmediata: regla de la potencia}
\small Ejercicio 7 \hfill Cálculo Integral -- TecNM
\ejercicio{$\displaystyle \int x^{-1/2}\,dx$}
\end{frame}

\begin{frame}{Integración inmediata: regla de la potencia}
\small Ejercicio 8 \hfill Cálculo Integral -- TecNM
\ejercicio{$\displaystyle \int x^{5/3}\,dx$}
\end{frame}

\begin{frame}{Polinomios y combinaciones lineales}
\small Ejercicio 1 \hfill Cálculo Integral -- TecNM
\ejercicio{$\displaystyle \int(3x^2-4x+7)\,dx$}
\end{frame}

\begin{frame}{Polinomios y combinaciones lineales}
\small Ejercicio 2 \hfill Cálculo Integral -- TecNM
\ejercicio{$\displaystyle \int(5x^4-2x^3+x-6)\,dx$}
\end{frame}

\begin{frame}{Polinomios y combinaciones lineales}
\small Ejercicio 3 \hfill Cálculo Integral -- TecNM
\ejercicio{$\displaystyle \int(2x^5+7x^2-9)\,dx$}
\end{frame}

\begin{frame}{Polinomios y combinaciones lineales}
\small Ejercicio 4 \hfill Cálculo Integral -- TecNM
\ejercicio{$\displaystyle \int(8-3x+4x^3)\,dx$}
\end{frame}

\begin{frame}{Polinomios y combinaciones lineales}
\small Ejercicio 5 \hfill Cálculo Integral -- TecNM
\ejercicio{$\displaystyle \int\left(\frac{x^4}{2}-\frac{3x^2}{4}+5x\right)dx$}
\end{frame}

\begin{frame}{Polinomios y combinaciones lineales}
\small Ejercicio 6 \hfill Cálculo Integral -- TecNM
\ejercicio{$\displaystyle \int(6x^7-5x^5+2x^2-1)\,dx$}
\end{frame}

\begin{frame}{Polinomios y combinaciones lineales}
\small Ejercicio 7 \hfill Cálculo Integral -- TecNM
\ejercicio{$\displaystyle \int\left(\frac{2}{3}x^6-\frac{5}{2}x^3+4\right)dx$}
\end{frame}

\begin{frame}{Polinomios y combinaciones lineales}
\small Ejercicio 8 \hfill Cálculo Integral -- TecNM
\ejercicio{$\displaystyle \int(7x^8-3x^4+x^2-8x+11)\,dx$}
\end{frame}

\begin{frame}{Radicales y potencias fraccionarias}
\small Ejercicio 1 \hfill Cálculo Integral -- TecNM
\ejercicio{$\displaystyle \int \sqrt{x}\,dx$}
\end{frame}

\begin{frame}{Radicales y potencias fraccionarias}
\small Ejercicio 2 \hfill Cálculo Integral -- TecNM
\ejercicio{$\displaystyle \int \sqrt[3]{x^2}\,dx$}
\end{frame}

\begin{frame}{Radicales y potencias fraccionarias}
\small Ejercicio 3 \hfill Cálculo Integral -- TecNM
\ejercicio{$\displaystyle \int \frac{1}{\sqrt{x}}\,dx$}
\end{frame}

\begin{frame}{Radicales y potencias fraccionarias}
\small Ejercicio 4 \hfill Cálculo Integral -- TecNM
\ejercicio{$\displaystyle \int \frac{1}{\sqrt[3]{x^2}}\,dx$}
\end{frame}

\begin{frame}{Radicales y potencias fraccionarias}
\small Ejercicio 5 \hfill Cálculo Integral -- TecNM
\ejercicio{$\displaystyle \int(2\sqrt{x}+3\sqrt[3]{x})\,dx$}
\end{frame}

\begin{frame}{Radicales y potencias fraccionarias}
\small Ejercicio 6 \hfill Cálculo Integral -- TecNM
\ejercicio{$\displaystyle \int\frac{x^2+\sqrt{x}}{x}\,dx$}
\end{frame}

\begin{frame}{Radicales y potencias fraccionarias}
\small Ejercicio 7 \hfill Cálculo Integral -- TecNM
\ejercicio{$\displaystyle \int\left(\sqrt[4]{x^3}+\frac{2}{\sqrt{x}}\right)dx$}
\end{frame}

\begin{frame}{Radicales y potencias fraccionarias}
\small Ejercicio 8 \hfill Cálculo Integral -- TecNM
\ejercicio{$\displaystyle \int\frac{x^3+2x}{\sqrt{x}}\,dx$}
\end{frame}

\begin{frame}{Simplificación algebraica antes de integrar}
\small Ejercicio 1 \hfill Cálculo Integral -- TecNM
\ejercicio{$\displaystyle \int\frac{x^3+2x^2-x}{x}\,dx$}
\end{frame}

\begin{frame}{Simplificación algebraica antes de integrar}
\small Ejercicio 2 \hfill Cálculo Integral -- TecNM
\ejercicio{$\displaystyle \int\frac{3x^4-2x^2+5x}{x^2}\,dx$}
\end{frame}

\begin{frame}{Simplificación algebraica antes de integrar}
\small Ejercicio 3 \hfill Cálculo Integral -- TecNM
\ejercicio{$\displaystyle \int\frac{(x+1)^2}{x}\,dx$}
\end{frame}

\begin{frame}{Simplificación algebraica antes de integrar}
\small Ejercicio 4 \hfill Cálculo Integral -- TecNM
\ejercicio{$\displaystyle \int\frac{x^4-1}{x^2}\,dx$}
\end{frame}

\begin{frame}{Simplificación algebraica antes de integrar}
\small Ejercicio 5 \hfill Cálculo Integral -- TecNM
\ejercicio{$\displaystyle \int\frac{2x^3+5x^2-3x}{x}\,dx$}
\end{frame}

\begin{frame}{Simplificación algebraica antes de integrar}
\small Ejercicio 6 \hfill Cálculo Integral -- TecNM
\ejercicio{$\displaystyle \int\frac{x^5+x^3+x}{x^2}\,dx$}
\end{frame}

\begin{frame}{Simplificación algebraica antes de integrar}
\small Ejercicio 7 \hfill Cálculo Integral -- TecNM
\ejercicio{$\displaystyle \int\frac{(x-2)(x+3)}{x}\,dx$}
\end{frame}

\begin{frame}{Simplificación algebraica antes de integrar}
\small Ejercicio 8 \hfill Cálculo Integral -- TecNM
\ejercicio{$\displaystyle \int\frac{x^6-2x^3+x}{x^2}\,dx$}
\end{frame}

\begin{frame}{La forma $1/x$ y logaritmos inmediatos}
\small Ejercicio 1 \hfill Cálculo Integral -- TecNM
\ejercicio{$\displaystyle \int\frac{1}{x}\,dx$}
\end{frame}

\begin{frame}{La forma $1/x$ y logaritmos inmediatos}
\small Ejercicio 2 \hfill Cálculo Integral -- TecNM
\ejercicio{$\displaystyle \int\frac{5}{x}\,dx$}
\end{frame}

\begin{frame}{La forma $1/x$ y logaritmos inmediatos}
\small Ejercicio 3 \hfill Cálculo Integral -- TecNM
\ejercicio{$\displaystyle \int\left(3x^2+\frac{2}{x}\right)dx$}
\end{frame}

\begin{frame}{La forma $1/x$ y logaritmos inmediatos}
\small Ejercicio 4 \hfill Cálculo Integral -- TecNM
\ejercicio{$\displaystyle \int\left(\frac{4}{x}-x^3\right)dx$}
\end{frame}

\begin{frame}{La forma $1/x$ y logaritmos inmediatos}
\small Ejercicio 5 \hfill Cálculo Integral -- TecNM
\ejercicio{$\displaystyle \int\frac{7x^2+3}{x}\,dx$}
\end{frame}

\begin{frame}{La forma $1/x$ y logaritmos inmediatos}
\small Ejercicio 6 \hfill Cálculo Integral -- TecNM
\ejercicio{$\displaystyle \int\frac{x^3-4}{x}\,dx$}
\end{frame}

\begin{frame}{La forma $1/x$ y logaritmos inmediatos}
\small Ejercicio 7 \hfill Cálculo Integral -- TecNM
\ejercicio{$\displaystyle \int\left(\sqrt{x}+\frac{6}{x}\right)dx$}
\end{frame}

\begin{frame}{La forma $1/x$ y logaritmos inmediatos}
\small Ejercicio 8 \hfill Cálculo Integral -- TecNM
\ejercicio{$\displaystyle \int\frac{2x^4-x+5}{x}\,dx$}
\end{frame}

\begin{frame}{Funciones exponenciales}
\small Ejercicio 1 \hfill Cálculo Integral -- TecNM
\ejercicio{$\displaystyle \int e^x\,dx$}
\end{frame}

\begin{frame}{Funciones exponenciales}
\small Ejercicio 2 \hfill Cálculo Integral -- TecNM
\ejercicio{$\displaystyle \int e^{3x}\,dx$}
\end{frame}

\begin{frame}{Funciones exponenciales}
\small Ejercicio 3 \hfill Cálculo Integral -- TecNM
\ejercicio{$\displaystyle \int e^{-2x}\,dx$}
\end{frame}

\begin{frame}{Funciones exponenciales}
\small Ejercicio 4 \hfill Cálculo Integral -- TecNM
\ejercicio{$\displaystyle \int 4e^{5x}\,dx$}
\end{frame}

\begin{frame}{Funciones exponenciales}
\small Ejercicio 5 \hfill Cálculo Integral -- TecNM
\ejercicio{$\displaystyle \int 2^x\,dx$}
\end{frame}

\begin{frame}{Funciones exponenciales}
\small Ejercicio 6 \hfill Cálculo Integral -- TecNM
\ejercicio{$\displaystyle \int 5^{3x}\,dx$}
\end{frame}

\begin{frame}{Funciones exponenciales}
\small Ejercicio 7 \hfill Cálculo Integral -- TecNM
\ejercicio{$\displaystyle \int 3^{-2x}\,dx$}
\end{frame}

\begin{frame}{Funciones exponenciales}
\small Ejercicio 8 \hfill Cálculo Integral -- TecNM
\ejercicio{$\displaystyle \int(2e^{4x}-3\,7^x)\,dx$}
\end{frame}

\begin{frame}{Trigonométricas inmediatas}
\small Ejercicio 1 \hfill Cálculo Integral -- TecNM
\ejercicio{$\displaystyle \int\sin x\,dx$}
\end{frame}

\begin{frame}{Trigonométricas inmediatas}
\small Ejercicio 2 \hfill Cálculo Integral -- TecNM
\ejercicio{$\displaystyle \int\cos x\,dx$}
\end{frame}

\begin{frame}{Trigonométricas inmediatas}
\small Ejercicio 3 \hfill Cálculo Integral -- TecNM
\ejercicio{$\displaystyle \int\sec^2x\,dx$}
\end{frame}

\begin{frame}{Trigonométricas inmediatas}
\small Ejercicio 4 \hfill Cálculo Integral -- TecNM
\ejercicio{$\displaystyle \int\csc^2x\,dx$}
\end{frame}

\begin{frame}{Trigonométricas inmediatas}
\small Ejercicio 5 \hfill Cálculo Integral -- TecNM
\ejercicio{$\displaystyle \int\sec x\tan x\,dx$}
\end{frame}

\begin{frame}{Trigonométricas inmediatas}
\small Ejercicio 6 \hfill Cálculo Integral -- TecNM
\ejercicio{$\displaystyle \int\csc x\cot x\,dx$}
\end{frame}

\begin{frame}{Trigonométricas inmediatas}
\small Ejercicio 7 \hfill Cálculo Integral -- TecNM
\ejercicio{$\displaystyle \int(3\cos x-2\sin x)\,dx$}
\end{frame}

\begin{frame}{Trigonométricas inmediatas}
\small Ejercicio 8 \hfill Cálculo Integral -- TecNM
\ejercicio{$\displaystyle \int(4\sec^2x+5\csc x\cot x)\,dx$}
\end{frame}

\begin{frame}{Trigonométricas con identidades elementales}
\small Ejercicio 1 \hfill Cálculo Integral -- TecNM
\ejercicio{$\displaystyle \int\sin^2x\,dx$}
\end{frame}

\begin{frame}{Trigonométricas con identidades elementales}
\small Ejercicio 2 \hfill Cálculo Integral -- TecNM
\ejercicio{$\displaystyle \int\cos^2x\,dx$}
\end{frame}

\begin{frame}{Trigonométricas con identidades elementales}
\small Ejercicio 3 \hfill Cálculo Integral -- TecNM
\ejercicio{$\displaystyle \int(1+\cos 2x)\,dx$}
\end{frame}

\begin{frame}{Trigonométricas con identidades elementales}
\small Ejercicio 4 \hfill Cálculo Integral -- TecNM
\ejercicio{$\displaystyle \int(1-\cos 2x)\,dx$}
\end{frame}

\begin{frame}{Trigonométricas con identidades elementales}
\small Ejercicio 5 \hfill Cálculo Integral -- TecNM
\ejercicio{$\displaystyle \int\sin x\cos x\,dx$}
\end{frame}

\begin{frame}{Trigonométricas con identidades elementales}
\small Ejercicio 6 \hfill Cálculo Integral -- TecNM
\ejercicio{$\displaystyle \int\sin(2x)\,dx$}
\end{frame}

\begin{frame}{Trigonométricas con identidades elementales}
\small Ejercicio 7 \hfill Cálculo Integral -- TecNM
\ejercicio{$\displaystyle \int\cos(5x)\,dx$}
\end{frame}

\begin{frame}{Trigonométricas con identidades elementales}
\small Ejercicio 8 \hfill Cálculo Integral -- TecNM
\ejercicio{$\displaystyle \int(3\sin^2x+2\cos^2x)\,dx$}
\end{frame}

\begin{frame}{Sustitución: potencias de una función}
\small Ejercicio 1 \hfill Cálculo Integral -- TecNM
\ejercicio{$\displaystyle \int 2x(x^2+1)^5\,dx$}
\end{frame}

\begin{frame}{Sustitución: potencias de una función}
\small Ejercicio 2 \hfill Cálculo Integral -- TecNM
\ejercicio{$\displaystyle \int 3x^2(x^3-4)^6\,dx$}
\end{frame}

\begin{frame}{Sustitución: potencias de una función}
\small Ejercicio 3 \hfill Cálculo Integral -- TecNM
\ejercicio{$\displaystyle \int(2x+1)(x^2+x+3)^4\,dx$}
\end{frame}

\begin{frame}{Sustitución: potencias de una función}
\small Ejercicio 4 \hfill Cálculo Integral -- TecNM
\ejercicio{$\displaystyle \int 4x^3(x^4+2)^7\,dx$}
\end{frame}

\begin{frame}{Sustitución: potencias de una función}
\small Ejercicio 5 \hfill Cálculo Integral -- TecNM
\ejercicio{$\displaystyle \int x(3x^2-1)^8\,dx$}
\end{frame}

\begin{frame}{Sustitución: potencias de una función}
\small Ejercicio 6 \hfill Cálculo Integral -- TecNM
\ejercicio{$\displaystyle \int(6x-5)(3x^2-5x+2)^3\,dx$}
\end{frame}

\begin{frame}{Sustitución: potencias de una función}
\small Ejercicio 7 \hfill Cálculo Integral -- TecNM
\ejercicio{$\displaystyle \int x^2(1-x^3)^5\,dx$}
\end{frame}

\begin{frame}{Sustitución: potencias de una función}
\small Ejercicio 8 \hfill Cálculo Integral -- TecNM
\ejercicio{$\displaystyle \int(5x^4+2)(x^5+2x-1)^6\,dx$}
\end{frame}

\begin{frame}{Sustitución con radicales}
\small Ejercicio 1 \hfill Cálculo Integral -- TecNM
\ejercicio{$\displaystyle \int\frac{x}{\sqrt{x^2+4}}\,dx$}
\end{frame}

\begin{frame}{Sustitución con radicales}
\small Ejercicio 2 \hfill Cálculo Integral -- TecNM
\ejercicio{$\displaystyle \int\frac{2x+1}{\sqrt{x^2+x+5}}\,dx$}
\end{frame}

\begin{frame}{Sustitución con radicales}
\small Ejercicio 3 \hfill Cálculo Integral -- TecNM
\ejercicio{$\displaystyle \int x^2\sqrt{x^3+1}\,dx$}
\end{frame}

\begin{frame}{Sustitución con radicales}
\small Ejercicio 4 \hfill Cálculo Integral -- TecNM
\ejercicio{$\displaystyle \int\frac{x^3}{\sqrt{x^4+7}}\,dx$}
\end{frame}

\begin{frame}{Sustitución con radicales}
\small Ejercicio 5 \hfill Cálculo Integral -- TecNM
\ejercicio{$\displaystyle \int(3x^2-2)\sqrt{x^3-2x+4}\,dx$}
\end{frame}

\begin{frame}{Sustitución con radicales}
\small Ejercicio 6 \hfill Cálculo Integral -- TecNM
\ejercicio{$\displaystyle \int\frac{4x}{\sqrt[3]{2x^2+1}}\,dx$}
\end{frame}

\begin{frame}{Sustitución con radicales}
\small Ejercicio 7 \hfill Cálculo Integral -- TecNM
\ejercicio{$\displaystyle \int x^2\sqrt[4]{x^3+2}\,dx$}
\end{frame}

\begin{frame}{Sustitución con radicales}
\small Ejercicio 8 \hfill Cálculo Integral -- TecNM
\ejercicio{$\displaystyle \int\frac{5x^4}{\sqrt{1+x^5}}\,dx$}
\end{frame}

\begin{frame}{Sustitución que conduce a logaritmos: $f'(x)/f(x)$}
\small Ejercicio 1 \hfill Cálculo Integral -- TecNM
\ejercicio{$\displaystyle \int\frac{2x}{x^2+3}\,dx$}
\end{frame}

\begin{frame}{Sustitución que conduce a logaritmos: $f'(x)/f(x)$}
\small Ejercicio 2 \hfill Cálculo Integral -- TecNM
\ejercicio{$\displaystyle \int\frac{3x^2}{x^3+5}\,dx$}
\end{frame}

\begin{frame}{Sustitución que conduce a logaritmos: $f'(x)/f(x)$}
\small Ejercicio 3 \hfill Cálculo Integral -- TecNM
\ejercicio{$\displaystyle \int\frac{2x+1}{x^2+x+4}\,dx$}
\end{frame}

\begin{frame}{Sustitución que conduce a logaritmos: $f'(x)/f(x)$}
\small Ejercicio 4 \hfill Cálculo Integral -- TecNM
\ejercicio{$\displaystyle \int\frac{4x^3}{x^4-7}\,dx$}
\end{frame}

\begin{frame}{Sustitución que conduce a logaritmos: $f'(x)/f(x)$}
\small Ejercicio 5 \hfill Cálculo Integral -- TecNM
\ejercicio{$\displaystyle \int\frac{6x-2}{3x^2-2x+9}\,dx$}
\end{frame}

\begin{frame}{Sustitución que conduce a logaritmos: $f'(x)/f(x)$}
\small Ejercicio 6 \hfill Cálculo Integral -- TecNM
\ejercicio{$\displaystyle \int\frac{5x^4+1}{x^5+x+2}\,dx$}
\end{frame}

\begin{frame}{Sustitución que conduce a logaritmos: $f'(x)/f(x)$}
\small Ejercicio 7 \hfill Cálculo Integral -- TecNM
\ejercicio{$\displaystyle \int\frac{\cos x}{2+\sin x}\,dx$}
\end{frame}

\begin{frame}{Sustitución que conduce a logaritmos: $f'(x)/f(x)$}
\small Ejercicio 8 \hfill Cálculo Integral -- TecNM
\ejercicio{$\displaystyle \int\frac{e^x}{1+e^x}\,dx$}
\end{frame}

\begin{frame}{Sustitución con exponenciales}
\small Ejercicio 1 \hfill Cálculo Integral -- TecNM
\ejercicio{$\displaystyle \int e^x(1+e^x)^4\,dx$}
\end{frame}

\begin{frame}{Sustitución con exponenciales}
\small Ejercicio 2 \hfill Cálculo Integral -- TecNM
\ejercicio{$\displaystyle \int e^{2x}\sqrt{1+e^{2x}}\,dx$}
\end{frame}

\begin{frame}{Sustitución con exponenciales}
\small Ejercicio 3 \hfill Cálculo Integral -- TecNM
\ejercicio{$\displaystyle \int\frac{e^x}{(3+e^x)^2}\,dx$}
\end{frame}

\begin{frame}{Sustitución con exponenciales}
\small Ejercicio 4 \hfill Cálculo Integral -- TecNM
\ejercicio{$\displaystyle \int e^{-x}(2+e^{-x})^5\,dx$}
\end{frame}

\begin{frame}{Sustitución con exponenciales}
\small Ejercicio 5 \hfill Cálculo Integral -- TecNM
\ejercicio{$\displaystyle \int\frac{3e^{3x}}{1+e^{3x}}\,dx$}
\end{frame}

\begin{frame}{Sustitución con exponenciales}
\small Ejercicio 6 \hfill Cálculo Integral -- TecNM
\ejercicio{$\displaystyle \int e^{4x}(1-e^{4x})^3\,dx$}
\end{frame}

\begin{frame}{Sustitución con exponenciales}
\small Ejercicio 7 \hfill Cálculo Integral -- TecNM
\ejercicio{$\displaystyle \int\frac{2^x}{1+2^x}\,dx$}
\end{frame}

\begin{frame}{Sustitución con exponenciales}
\small Ejercicio 8 \hfill Cálculo Integral -- TecNM
\ejercicio{$\displaystyle \int 5^x(2+5^x)^6\,dx$}
\end{frame}

\begin{frame}{Sustitución con funciones trigonométricas}
\small Ejercicio 1 \hfill Cálculo Integral -- TecNM
\ejercicio{$\displaystyle \int\sin x\cos^4x\,dx$}
\end{frame}

\begin{frame}{Sustitución con funciones trigonométricas}
\small Ejercicio 2 \hfill Cálculo Integral -- TecNM
\ejercicio{$\displaystyle \int\cos x\sin^5x\,dx$}
\end{frame}

\begin{frame}{Sustitución con funciones trigonométricas}
\small Ejercicio 3 \hfill Cálculo Integral -- TecNM
\ejercicio{$\displaystyle \int\sec^2x(1+\tan x)^3\,dx$}
\end{frame}

\begin{frame}{Sustitución con funciones trigonométricas}
\small Ejercicio 4 \hfill Cálculo Integral -- TecNM
\ejercicio{$\displaystyle \int\csc^2x(2-\cot x)^4\,dx$}
\end{frame}

\begin{frame}{Sustitución con funciones trigonométricas}
\small Ejercicio 5 \hfill Cálculo Integral -- TecNM
\ejercicio{$\displaystyle \int\sec x\tan x(3+\sec x)^5\,dx$}
\end{frame}

\begin{frame}{Sustitución con funciones trigonométricas}
\small Ejercicio 6 \hfill Cálculo Integral -- TecNM
\ejercicio{$\displaystyle \int\csc x\cot x(1+\csc x)^2\,dx$}
\end{frame}

\begin{frame}{Sustitución con funciones trigonométricas}
\small Ejercicio 7 \hfill Cálculo Integral -- TecNM
\ejercicio{$\displaystyle \int\frac{\cos x}{1+\sin x}\,dx$}
\end{frame}

\begin{frame}{Sustitución con funciones trigonométricas}
\small Ejercicio 8 \hfill Cálculo Integral -- TecNM
\ejercicio{$\displaystyle \int\frac{\sec^2x}{2+\tan x}\,dx$}
\end{frame}

\begin{frame}{Integrales que conducen a funciones trigonométricas inversas}
\small Ejercicio 1 \hfill Cálculo Integral -- TecNM
\ejercicio{$\displaystyle \int\frac{1}{1+x^2}\,dx$}
\end{frame}

\begin{frame}{Integrales que conducen a funciones trigonométricas inversas}
\small Ejercicio 2 \hfill Cálculo Integral -- TecNM
\ejercicio{$\displaystyle \int\frac{1}{\sqrt{1-x^2}}\,dx$}
\end{frame}

\begin{frame}{Integrales que conducen a funciones trigonométricas inversas}
\small Ejercicio 3 \hfill Cálculo Integral -- TecNM
\ejercicio{$\displaystyle \int\frac{1}{4+x^2}\,dx$}
\end{frame}

\begin{frame}{Integrales que conducen a funciones trigonométricas inversas}
\small Ejercicio 4 \hfill Cálculo Integral -- TecNM
\ejercicio{$\displaystyle \int\frac{1}{\sqrt{9-x^2}}\,dx$}
\end{frame}

\begin{frame}{Integrales que conducen a funciones trigonométricas inversas}
\small Ejercicio 5 \hfill Cálculo Integral -- TecNM
\ejercicio{$\displaystyle \int\frac{1}{1+4x^2}\,dx$}
\end{frame}

\begin{frame}{Integrales que conducen a funciones trigonométricas inversas}
\small Ejercicio 6 \hfill Cálculo Integral -- TecNM
\ejercicio{$\displaystyle \int\frac{1}{\sqrt{16-9x^2}}\,dx$}
\end{frame}

\begin{frame}{Integrales que conducen a funciones trigonométricas inversas}
\small Ejercicio 7 \hfill Cálculo Integral -- TecNM
\ejercicio{$\displaystyle \int\frac{3}{1+9x^2}\,dx$}
\end{frame}

\begin{frame}{Integrales que conducen a funciones trigonométricas inversas}
\small Ejercicio 8 \hfill Cálculo Integral -- TecNM
\ejercicio{$\displaystyle \int\frac{2}{\sqrt{25-4x^2}}\,dx$}
\end{frame}

\begin{frame}{Integrales mixtas: identificar el método}
\small Ejercicio 1 \hfill Cálculo Integral -- TecNM
\ejercicio{$\displaystyle \int(4x^3-2x+7)\,dx$}
\end{frame}

\begin{frame}{Integrales mixtas: identificar el método}
\small Ejercicio 2 \hfill Cálculo Integral -- TecNM
\ejercicio{$\displaystyle \int\frac{x^2+3x}{x}\,dx$}
\end{frame}

\begin{frame}{Integrales mixtas: identificar el método}
\small Ejercicio 3 \hfill Cálculo Integral -- TecNM
\ejercicio{$\displaystyle \int\frac{x}{x^2+9}\,dx$}
\end{frame}

\begin{frame}{Integrales mixtas: identificar el método}
\small Ejercicio 4 \hfill Cálculo Integral -- TecNM
\ejercicio{$\displaystyle \int x^2(x^3+4)^7\,dx$}
\end{frame}

\begin{frame}{Integrales mixtas: identificar el método}
\small Ejercicio 5 \hfill Cálculo Integral -- TecNM
\ejercicio{$\displaystyle \int\frac{x}{\sqrt{x^2+1}}\,dx$}
\end{frame}

\begin{frame}{Integrales mixtas: identificar el método}
\small Ejercicio 6 \hfill Cálculo Integral -- TecNM
\ejercicio{$\displaystyle \int e^{2x}(1+e^{2x})^3\,dx$}
\end{frame}

\begin{frame}{Integrales mixtas: identificar el método}
\small Ejercicio 7 \hfill Cálculo Integral -- TecNM
\ejercicio{$\displaystyle \int\cos x(2+\sin x)^5\,dx$}
\end{frame}

\begin{frame}{Integrales mixtas: identificar el método}
\small Ejercicio 8 \hfill Cálculo Integral -- TecNM
\ejercicio{$\displaystyle \int\frac{\sec^2x}{1+\tan x}\,dx$}
\end{frame}

\begin{frame}{Integrales mixtas: identificar el método}
\small Ejercicio 9 \hfill Cálculo Integral -- TecNM
\ejercicio{$\displaystyle \int\cos^2x\,dx$}
\end{frame}

\begin{frame}{Integrales mixtas: identificar el método}
\small Ejercicio 10 \hfill Cálculo Integral -- TecNM
\ejercicio{$\displaystyle \int\frac{1}{9+x^2}\,dx$}
\end{frame}

\begin{frame}{Integrales mixtas: identificar el método}
\small Ejercicio 11 \hfill Cálculo Integral -- TecNM
\ejercicio{$\displaystyle \int\frac{1}{\sqrt{25-x^2}}\,dx$}
\end{frame}

\begin{frame}{Integrales mixtas: identificar el método}
\small Ejercicio 12 \hfill Cálculo Integral -- TecNM
\ejercicio{$\displaystyle \int\left(x^{3/2}+x^{-1/2}\right)dx$}
\end{frame}

\begin{frame}{Integrales mixtas: identificar el método}
\small Ejercicio 13 \hfill Cálculo Integral -- TecNM
\ejercicio{$\displaystyle \int 7^{2x}\,dx$}
\end{frame}

\begin{frame}{Integrales mixtas: identificar el método}
\small Ejercicio 14 \hfill Cálculo Integral -- TecNM
\ejercicio{$\displaystyle \int\frac{4x^3}{2+x^4}\,dx$}
\end{frame}

\begin{frame}{Integrales mixtas: identificar el método}
\small Ejercicio 15 \hfill Cálculo Integral -- TecNM
\ejercicio{$\displaystyle \int\sin x\cos^6x\,dx$}
\end{frame}

\begin{frame}{Integrales mixtas: identificar el método}
\small Ejercicio 16 \hfill Cálculo Integral -- TecNM
\ejercicio{$\displaystyle \int\frac{e^x}{4+e^x}\,dx$}
\end{frame}

\begin{frame}{Retos de cierre}
\small Ejercicio 1 \hfill Cálculo Integral -- TecNM
\ejercicio{$\displaystyle \int\frac{x^5+2x^3-x}{x^2}\,dx$}
\end{frame}

\begin{frame}{Retos de cierre}
\small Ejercicio 2 \hfill Cálculo Integral -- TecNM
\ejercicio{$\displaystyle \int\frac{x^2}{\sqrt{1+x^3}}\,dx$}
\end{frame}

\begin{frame}{Retos de cierre}
\small Ejercicio 3 \hfill Cálculo Integral -- TecNM
\ejercicio{$\displaystyle \int\frac{2x-3}{x^2-3x+10}\,dx$}
\end{frame}

\begin{frame}{Retos de cierre}
\small Ejercicio 4 \hfill Cálculo Integral -- TecNM
\ejercicio{$\displaystyle \int e^{-3x}(1+e^{-3x})^7\,dx$}
\end{frame}

\begin{frame}{Retos de cierre}
\small Ejercicio 5 \hfill Cálculo Integral -- TecNM
\ejercicio{$\displaystyle \int\frac{\cos(2x)}{3+\sin(2x)}\,dx$}
\end{frame}

\begin{frame}{Retos de cierre}
\small Ejercicio 6 \hfill Cálculo Integral -- TecNM
\ejercicio{$\displaystyle \int\frac{1}{16+25x^2}\,dx$}
\end{frame}

\begin{frame}{Retos de cierre}
\small Ejercicio 7 \hfill Cálculo Integral -- TecNM
\ejercicio{$\displaystyle \int\frac{1}{\sqrt{36-4x^2}}\,dx$}
\end{frame}

\begin{frame}{Retos de cierre}
\small Ejercicio 8 \hfill Cálculo Integral -- TecNM
\ejercicio{$\displaystyle \int\left(2\sin^2x+3\cos^2x\right)dx$}
\end{frame}


%=========================================================
% Potencias de seno
%=========================================================
\begin{frame}{Potencias de seno}
\small Ejercicio 1 \hfill Cálculo Integral -- TecNM
\ejercicio{$\displaystyle \int \sin^3 x\,dx$}
\end{frame}
\begin{frame}{Potencias de seno}
\small Ejercicio 2 \hfill Cálculo Integral -- TecNM
\ejercicio{$\displaystyle \int \sin^5 x\,dx$}
\end{frame}
\begin{frame}{Potencias de seno}
\small Ejercicio 3 \hfill Cálculo Integral -- TecNM
\ejercicio{$\displaystyle \int \sin^7 x\,dx$}
\end{frame}
\begin{frame}{Potencias de seno}
\small Ejercicio 4 \hfill Cálculo Integral -- TecNM
\ejercicio{$\displaystyle \int \sin^2 x\,dx$}
\end{frame}
\begin{frame}{Potencias de seno}
\small Ejercicio 5 \hfill Cálculo Integral -- TecNM
\ejercicio{$\displaystyle \int \sin^4 x\,dx$}
\end{frame}
\begin{frame}{Potencias de seno}
\small Ejercicio 6 \hfill Cálculo Integral -- TecNM
\ejercicio{$\displaystyle \int \sin^6 x\,dx$}
\end{frame}
\begin{frame}{Potencias de seno}
\small Ejercicio 7 \hfill Cálculo Integral -- TecNM
\ejercicio{$\displaystyle \int \sin^3(2x)\,dx$}
\end{frame}
\begin{frame}{Potencias de seno}
\small Ejercicio 8 \hfill Cálculo Integral -- TecNM
\ejercicio{$\displaystyle \int \sin^4(3x)\,dx$}
\end{frame}

%=========================================================
% Potencias de coseno
%=========================================================
\begin{frame}{Potencias de coseno}
\small Ejercicio 1 \hfill Cálculo Integral -- TecNM
\ejercicio{$\displaystyle \int \cos^3 x\,dx$}
\end{frame}
\begin{frame}{Potencias de coseno}
\small Ejercicio 2 \hfill Cálculo Integral -- TecNM
\ejercicio{$\displaystyle \int \cos^5 x\,dx$}
\end{frame}
\begin{frame}{Potencias de coseno}
\small Ejercicio 3 \hfill Cálculo Integral -- TecNM
\ejercicio{$\displaystyle \int \cos^7 x\,dx$}
\end{frame}
\begin{frame}{Potencias de coseno}
\small Ejercicio 4 \hfill Cálculo Integral -- TecNM
\ejercicio{$\displaystyle \int \cos^2 x\,dx$}
\end{frame}
\begin{frame}{Potencias de coseno}
\small Ejercicio 5 \hfill Cálculo Integral -- TecNM
\ejercicio{$\displaystyle \int \cos^4 x\,dx$}
\end{frame}
\begin{frame}{Potencias de coseno}
\small Ejercicio 6 \hfill Cálculo Integral -- TecNM
\ejercicio{$\displaystyle \int \cos^6 x\,dx$}
\end{frame}
\begin{frame}{Potencias de coseno}
\small Ejercicio 7 \hfill Cálculo Integral -- TecNM
\ejercicio{$\displaystyle \int \cos^3(2x)\,dx$}
\end{frame}
\begin{frame}{Potencias de coseno}
\small Ejercicio 8 \hfill Cálculo Integral -- TecNM
\ejercicio{$\displaystyle \int \cos^4(3x)\,dx$}
\end{frame}

%=========================================================
% Productos de potencias seno-coseno
%=========================================================
\begin{frame}{Productos de potencias de seno y coseno}
\small Ejercicio 1 \hfill Cálculo Integral -- TecNM
\ejercicio{$\displaystyle \int \sin^3 x\cos^2 x\,dx$}
\end{frame}
\begin{frame}{Productos de potencias de seno y coseno}
\small Ejercicio 2 \hfill Cálculo Integral -- TecNM
\ejercicio{$\displaystyle \int \sin^2 x\cos^3 x\,dx$}
\end{frame}
\begin{frame}{Productos de potencias de seno y coseno}
\small Ejercicio 3 \hfill Cálculo Integral -- TecNM
\ejercicio{$\displaystyle \int \sin^5 x\cos^4 x\,dx$}
\end{frame}
\begin{frame}{Productos de potencias de seno y coseno}
\small Ejercicio 4 \hfill Cálculo Integral -- TecNM
\ejercicio{$\displaystyle \int \sin^4 x\cos^5 x\,dx$}
\end{frame}
\begin{frame}{Productos de potencias de seno y coseno}
\small Ejercicio 5 \hfill Cálculo Integral -- TecNM
\ejercicio{$\displaystyle \int \sin^2 x\cos^2 x\,dx$}
\end{frame}
\begin{frame}{Productos de potencias de seno y coseno}
\small Ejercicio 6 \hfill Cálculo Integral -- TecNM
\ejercicio{$\displaystyle \int \sin^4 x\cos^2 x\,dx$}
\end{frame}
\begin{frame}{Productos de potencias de seno y coseno}
\small Ejercicio 7 \hfill Cálculo Integral -- TecNM
\ejercicio{$\displaystyle \int \sin^2 x\cos^4 x\,dx$}
\end{frame}
\begin{frame}{Productos de potencias de seno y coseno}
\small Ejercicio 8 \hfill Cálculo Integral -- TecNM
\ejercicio{$\displaystyle \int \sin^4 x\cos^4 x\,dx$}
\end{frame}

%=========================================================
% Sustitución trigonométrica
%=========================================================
\begin{frame}{Sustitución trigonométrica}
\small Ejercicio 1 \hfill Cálculo Integral -- TecNM
\ejercicio{$\displaystyle \int \frac{dx}{\sqrt{9-x^2}}$}
\end{frame}
\begin{frame}{Sustitución trigonométrica}
\small Ejercicio 2 \hfill Cálculo Integral -- TecNM
\ejercicio{$\displaystyle \int \sqrt{16-x^2}\,dx$}
\end{frame}
\begin{frame}{Sustitución trigonométrica}
\small Ejercicio 3 \hfill Cálculo Integral -- TecNM
\ejercicio{$\displaystyle \int \frac{x^2}{\sqrt{25-x^2}}\,dx$}
\end{frame}
\begin{frame}{Sustitución trigonométrica}
\small Ejercicio 4 \hfill Cálculo Integral -- TecNM
\ejercicio{$\displaystyle \int \frac{dx}{\sqrt{x^2+9}}$}
\end{frame}
\begin{frame}{Sustitución trigonométrica}
\small Ejercicio 5 \hfill Cálculo Integral -- TecNM
\ejercicio{$\displaystyle \int \sqrt{x^2+4}\,dx$}
\end{frame}
\begin{frame}{Sustitución trigonométrica}
\small Ejercicio 6 \hfill Cálculo Integral -- TecNM
\ejercicio{$\displaystyle \int \frac{x^2}{\sqrt{x^2+16}}\,dx$}
\end{frame}
\begin{frame}{Sustitución trigonométrica}
\small Ejercicio 7 \hfill Cálculo Integral -- TecNM
\ejercicio{$\displaystyle \int \frac{dx}{\sqrt{x^2-9}}$}
\end{frame}
\begin{frame}{Sustitución trigonométrica}
\small Ejercicio 8 \hfill Cálculo Integral -- TecNM
\ejercicio{$\displaystyle \int \frac{\sqrt{x^2-25}}{x}\,dx$}
\end{frame}

%=========================================================
% Funciones trigonométricas inversas
%=========================================================
\begin{frame}{Integrales que conducen a $\arcsin$}
\small Ejercicio 1 \hfill Cálculo Integral -- TecNM
\ejercicio{$\displaystyle \int \frac{dx}{\sqrt{1-x^2}}$}
\end{frame}
\begin{frame}{Integrales que conducen a $\arcsin$}
\small Ejercicio 2 \hfill Cálculo Integral -- TecNM
\ejercicio{$\displaystyle \int \frac{dx}{\sqrt{4-x^2}}$}
\end{frame}
\begin{frame}{Integrales que conducen a $\arcsin$}
\small Ejercicio 3 \hfill Cálculo Integral -- TecNM
\ejercicio{$\displaystyle \int \frac{dx}{\sqrt{25-4x^2}}$}
\end{frame}
\begin{frame}{Integrales que conducen a $\arcsin$}
\small Ejercicio 4 \hfill Cálculo Integral -- TecNM
\ejercicio{$\displaystyle \int \frac{3\,dx}{\sqrt{16-9x^2}}$}
\end{frame}
\begin{frame}{Integrales que conducen a $\arctan$}
\small Ejercicio 5 \hfill Cálculo Integral -- TecNM
\ejercicio{$\displaystyle \int \frac{dx}{1+x^2}$}
\end{frame}
\begin{frame}{Integrales que conducen a $\arctan$}
\small Ejercicio 6 \hfill Cálculo Integral -- TecNM
\ejercicio{$\displaystyle \int \frac{dx}{9+x^2}$}
\end{frame}
\begin{frame}{Integrales que conducen a $\arctan$}
\small Ejercicio 7 \hfill Cálculo Integral -- TecNM
\ejercicio{$\displaystyle \int \frac{dx}{4+9x^2}$}
\end{frame}
\begin{frame}{Integrales que conducen a $\arctan$}
\small Ejercicio 8 \hfill Cálculo Integral -- TecNM
\ejercicio{$\displaystyle \int \frac{5\,dx}{25+4x^2}$}
\end{frame}

\end{document}
